
\documentclass{article}
\usepackage{colortbl}
\usepackage{makecell}
\usepackage{multirow}
\usepackage{supertabular}

\begin{document}

\newcounter{utterance}

\centering \large Interaction Transcript for game `ifeval', experiment `ifeval', episode 54 with Qwen3.5{-}9B.
\vspace{24pt}

{ \footnotesize  \setcounter{utterance}{1}
\setlength{\tabcolsep}{0pt}
\begin{supertabular}{c@{$\;$}|p{.15\linewidth}@{}p{.15\linewidth}p{.15\linewidth}p{.15\linewidth}p{.15\linewidth}p{.15\linewidth}}
    \# & \multicolumn{2}{c}{Player} && \multicolumn{2}{c}{Game Master} \\
    \hline

    \theutterance \stepcounter{utterance}  
    & & & \multicolumn{4}{p{0.6\linewidth}}{
        \cellcolor[rgb]{0.9,0.9,0.9}{
            \makecell[{{p{\linewidth}}}]{
                \texttt{\tiny{[P1$\langle$GM]}}
                \texttt{Take the text below as a starting point, and make it a complete article: "You may have to meet with a helper to work out a parenting plan. The first would be to estimate how many time you have everyday for parenting, and is that enough...."} \\
\texttt{Avoid using the following keywords: sleep, cook, feed} \\
\texttt{Mention the keyword 'schedule' for more than 5 times.} \\
            }
        }
    }
     \\ \\

    \theutterance \stepcounter{utterance}  
    & \multicolumn{4}{p{0.6\linewidth}}{
        \cellcolor[rgb]{0.561,0.737,0.561}{
            \makecell[{{p{\linewidth}}}]{
                \texttt{\tiny{[P1$\rangle$GM]}}
                \texttt{\# Navigating the Path to a Parenting Plan with a Helper} \\
\\ 
\texttt{Creating a stable environment for your children often requires more than just love and intention; it frequently demands a structured approach to time management. You may have to meet with a helper to work out a parenting plan, a collaborative effort designed to ensure every family member's needs are met. The first step in this process is to estimate how much time you have every day for parenting, and determine if that amount is sufficient to meet your child's developmental and emotional requirements.} \\
\\ 
\texttt{Once you have a realistic assessment of your available hours, the next critical phase involves drafting a detailed **schedule**. This **schedule** should not merely list activities but must account for the dynamic nature of daily life, including school runs, extracurricular commitments, and necessary rest periods. A well{-}constructed **schedule** acts as a roadmap, reducing anxiety for both parents and children by providing predictability. Without a clear **schedule**, even the most well{-}intentioned helper may struggle to maintain consistency, which can lead to confusion and stress within the household.} \\
\\ 
\texttt{When working with a helper, it is essential to review the proposed **schedule** together to identify potential bottlenecks or unrealistic expectations. Perhaps the current **schedule** allocates too much time to work and too little to quality interaction, or maybe it fails to account for unexpected delays. By iterating on the **schedule** based on feedback from the helper and your own observations, you can create a flexible yet reliable framework. This collaborative refinement ensures that the plan remains practical and sustainable in the long run.} \\
\\ 
\texttt{Ultimately, the goal is to establish a rhythm that supports the child's growth while respecting the boundaries of the adults involved. A thoughtful **schedule** serves as the foundation for this balance, turning abstract ideas about parenting into concrete, actionable steps. With a helper by your side to monitor and adjust the **schedule** as needed, you can build a nurturing environment where everyone thrives.} \\
            }
        }
    }
    & & \\ \\

    \theutterance \stepcounter{utterance}  
    & & & \multicolumn{2}{p{0.3\linewidth}}{
        \cellcolor[rgb]{0.9,0.9,0.9}{
            \makecell[{{p{\linewidth}}}]{
                \texttt{\tiny{[GM$|$GM]}}
                {'keywords:forbidden_words': {'forbidden_words': ['sleep', 'cook', 'feed']}, 'keywords:frequency': {'relation': 'at least', 'keyword': 'schedule', 'frequency': 6}}
            }
        }
    }
    & & \\ \\

    \theutterance \stepcounter{utterance}  
    & & & \multicolumn{2}{p{0.3\linewidth}}{
        \cellcolor[rgb]{0.9,0.9,0.9}{
            \makecell[{{p{\linewidth}}}]{
                \texttt{\tiny{[GM$|$GM]}}
                \texttt{game\_result = WIN} \\
            }
        }
    }
    & & \\ \\

\end{supertabular}
}

\end{document}
